% !TeX root = ../main.tex
\chapter{Virtual Vehicle and Digital-Twin Training Workflow}
\label{ch:system}
\section{Virtual Counterpart and Training Scope}
The virtual counterpart represents the vehicle geometry, differential-drive
motion, LiDAR configuration, command labels and task environment needed for
training. Shared parameters make these assumptions explicit and reproducible.
For example, wheel radius and separation determine the kinematic conversion,
while scan layout and distance units determine the student's input contract.
This coupling is the practical reason for using a vehicle-specific virtual
model instead of generating arbitrary range vectors with invented labels.

The role of this counterpart is to support interaction that would otherwise
require repeated driving of the physical vehicle. It can reset starts, expose
map geometry to the teacher and collect sensor/action pairs over long runs.
These capabilities support the intended digital-twin workflow. They do not
show that simulated friction, range noise or motor response have been fitted
to physical measurements. The present model is a parameterised training
representation, with physical fidelity and synchronisation still to be assessed.

The legacy physical firmware also uses different steering and timing conventions.
Consequently, the common simulation interface is not yet a demonstrated
end-to-end physical interface. The next integration stage must reconcile
those conventions and validate the virtual model against the actual vehicle.

\section{Training and Deployment Architecture}
The implementation combines ROS~2 communication~\cite{macenski2022ros2}
with Gazebo Classic physics, using the simulation framework introduced by
Koenig and Howard~\cite{koenig2004gazebo}.
A shared vehicle configuration defines geometry, sensor settings and steering
labels. Teacher training, demonstration collection and student evaluation use
these definitions, while saved run configurations provide the historical
settings for individual experiments. Current source defaults are not used as
substitutes for the configuration stored with an older model.

The learning sequence is shown in Figure~\ref{fig:pipeline}. During teacher
operation, map geometry and vehicle state enter the learned neural policy,
which selects an action. Collection
also records LiDAR, but LiDAR does not become an input to that privileged
policy. During student evaluation, the policy receives LiDAR only; geometry
and state remain available to the evaluator for scoring.

\begin{figure}[htbp]
\centering
\fbox{\parbox{0.89\textwidth}{\centering
\textbf{Virtual training on the host}\\[5pt]
Vehicle, sensor and task model $\longrightarrow$ simulated interaction\\[5pt]
Reward feedback $\longrightarrow$ PPO training of an MLP teacher\\[5pt]
Teacher commands + paired LiDAR $\longrightarrow$ labelled dataset\\[5pt]
Supervised fitting $\longrightarrow$ random forest / single tree\\[5pt]
Student closed-loop evaluation + C++ agreement checks\\[8pt]
$\downarrow$\\[5pt]
\textbf{Intended onboard inference}\\[5pt]
Physical LiDAR $\longrightarrow$ exported student $\longrightarrow$ motors\\[3pt]
Physical calibration and board validation remain pending}}
\caption[Teacher-to-student information flow]{Implemented information flow. Geometry used by the evaluator does
not enter the student's feature vector. The lower stage is the intended
physical integration path, not a completed experiment.}
\label{fig:pipeline}
\end{figure}

This separation avoids evaluating a purported LiDAR controller that silently
uses ground-truth distances. It also enables geometric clearance to be
computed independently of sensor corruption. A missing scan can affect the
policy without simultaneously hiding an obstacle from the scoring system.

\section{Differential-Drive Action Interface}
Let $a\in\{0,\ldots,14\}$ be a steering label, with $a=7$ denoting straight
motion. The implemented targets are
\begin{align}
 s(a)&=\frac{7-a}{7}, & \Delta v(a)&=\rho v_0 s(a),\\
 v_L(a)&=v_0-\Delta v(a), & v_R(a)&=v_0+\Delta v(a),
\end{align}
where $v_0=0.2\,\mathrm{m\,s^{-1}}$ and $\rho=0.6$. With wheel separation
$b=0.11\,\mathrm{m}$, the ideal differential-drive relation~\cite[Chapter~13]{lynch2017modern} gives
\begin{equation}
 v=\frac{v_R+v_L}{2}=v_0,\qquad
 \omega=\frac{v_R-v_L}{b}.
\end{equation}
Wheel angular targets are obtained by division by the configured radius
$r=0.0325\,\mathrm{m}$. At the strongest steering commands the linear wheel
targets are $0.08$ and $0.32\,\mathrm{m\,s^{-1}}$. Thus all ordinary steering
actions continue forward. An ideal maximum-curvature radius computed from
these targets is about $0.092\,\mathrm{m}$; actual collision-free passage also
depends on body dimensions, wall shape and the finite control interval.

The stop command is a separate interface value, 200, and is not a learned
steering class. The model therefore selects curvature at an approximately
fixed centre speed, rather than jointly optimising forward speed and turning.
This restriction limits the conclusions to roads that can be traversed with
forward-only control.

\section{LiDAR Representation and Missing Values}
Raw scan angles are first converted to integer degrees. For an angle $\theta$
in degrees, the feature index is
\begin{equation}
 i=\left\lfloor\frac{200\lfloor\theta\rfloor}{360}\right\rfloor,
 \qquad 0\leq\theta<360.
\end{equation}
Distances are represented in integer millimetres. Values outside the accepted
positive range below 10,000~mm are represented as invalid rather than extended
to arbitrary large measurements. The physical sensor range in the simulator
and the feature representation limit are distinct settings.

The active indices are $0$--$49$ and $151$--$199$. The remaining 101 indices
are fixed zero, leaving 99 active angular bins. Zero in an active bin denotes
a missing or unusable reading. A model trained on these inputs must therefore
interpret both angle and missingness consistently. Removing all rows containing
a zero would incorrectly discard valid scans because the inactive sector is
always zero.

Table~\ref{tab:platform} gives the settings used for the reported collections
and student evaluations. The 1~cm noise standard deviation is a simulation
assumption, not a measured calibration error. Earlier teacher configurations
recorded a different raw ray count, but the teacher did not use LiDAR; the
student dataset and evaluated sensor configuration use 364 rays per revolution.

\begin{table}[htbp]
\centering\small
\caption{Principal simulation settings for the student experiments.}
\label{tab:platform}
\begin{tabularx}{\textwidth}{Xl}
\toprule Parameter & Value \\\midrule
Target forward speed & $0.2\,\mathrm{m\,s^{-1}}$ \\
Wheel radius / separation & $0.0325/0.11\,\mathrm{m}$ \\
Steering labels / straight label & $0$--$14$ / $7$ \\
Decision interval / physics step & $0.2/0.001\,\mathrm{s}$ \\
LiDAR scan rate / raw rays & $5.5\,\mathrm{Hz}$ / $364$ \\
LiDAR physical range & $0.15$--$12\,\mathrm{m}$ \\
Feature distance upper limit & $10\,\mathrm{m}$, exclusive \\
Range-noise standard deviation & $0.01\,\mathrm{m}$ \\
Input columns / active bins & $200/99$ \\
Conservative body radius for clearance & $0.085\,\mathrm{m}$ \\
\bottomrule
\end{tabularx}
\end{table}

\section{Maps and Experimental Roles}
Four static ring roads are used: \mapname{train_world3},
\mapname{rounded_rectangle_001}, \mapname{rounded_rectangle_003} and
\mapname{ellipse_002}. For compact tables, these are denoted T3, R1, R3 and
E2, respectively. Wall geometry is represented by static obstacles, and
same-stem map metadata provide spawn poses and workspace bounds.

The selected teacher was trained and qualified on T3. It subsequently supplied
demonstrations on all four maps. Student fitting uses T3, R1 and R3; E2 is
withheld from fitting but used for validation and selection. The distinction
between teacher training, teacher collection, student fitting and student
selection prevents an ambiguous claim that a map is simply ``unseen''.
Figure~\ref{fig:collisionmaps} later shows the four geometries with recorded
collision endpoints.

\section{Timing and Evaluation Boundaries}
The formal closed-loop evaluator pauses Gazebo between controlled steps and
uses a nominal 0.2~s simulation decision interval. This supports repeatable
state/action records, but does not demonstrate wall-clock real-time operation.
The GUI driver operates asynchronously and is not the source of the formal
comparison matrices.

Scoring uses a conservative body disc and map-derived obstacle distance.
Clearance is sampled at decision steps, so its minimum is a sampled geometric
measure. Contact-based collision detection is a different signal. A reported
zero clearance can therefore occur without a recorded collision; neither
quantity should silently replace the other.

The intended hardware implementation must eventually reproduce the input and
action contract. The legacy firmware uses different steering and timing
conventions, so no physical equivalence is assumed here. All numerical results
in the following chapters concern the saved simulation configurations.
